\chapter{Triển khai công nghệ và Kết quả}

\section{Triển khai công nghệ}

Để xây dựng hệ thống Multi-Agent trợ lý y tế ảo với khả năng phân tích triệu chứng đa phương thức và cung cấp thông tin y khoa đáng tin cậy, đề tài lựa chọn tích hợp các công nghệ hiện đại, tối ưu cho từng thành phần của workflow AI phức tạp. Việc kết hợp các nền tảng và mô hình dưới đây giúp đảm bảo hiệu quả, khả năng mở rộng và độ tin cậy của hệ thống:

\begin{itemize}
    \item \textbf{LangGraph:} Framework mã nguồn mở cho phép xây dựng và điều phối các workflow AI phức tạp theo mô hình đồ thị. LangGraph hỗ trợ quản lý trạng thái, điều hướng luồng xử lý, tích hợp nhiều tác nhân AI và kiểm soát linh hoạt quá trình xử lý, phù hợp với các ứng dụng đa tác nhân và workflow có nhiều bước rẽ nhánh, lặp lại.
    
    \item \textbf{E5 Multilingual Large Instruct:} Mô hình embedding đa ngôn ngữ mạnh mẽ, hỗ trợ hơn 100 ngôn ngữ, được tối ưu cho các tác vụ tìm kiếm ngữ nghĩa và các bài toán nhúng văn bản phức tạp. Mô hình này tạo ra vector biểu diễn ngữ nghĩa chất lượng cao, phù hợp cho các hệ thống truy xuất thông tin đa ngôn ngữ, đặc biệt là tiếng Việt.
    
    \item \textbf{Qdrant Vector Database:} Cơ sở dữ liệu vector chuyên dụng, tối ưu cho lưu trữ và tìm kiếm dữ liệu vector hiệu suất cao. Qdrant hỗ trợ tìm kiếm theo độ tương đồng ngữ nghĩa, tích hợp các thuật toán tìm kiếm gần đúng (ANN) và cho phép kết hợp truy vấn vector với các bộ lọc dữ liệu bổ sung, rất phù hợp cho ứng dụng AI cần truy xuất thông tin chính xác, nhanh chóng.
    
    \item \textbf{Gemini 2.0 Flash:} Mô hình ngôn ngữ lớn (LLM) thế hệ mới, hỗ trợ xử lý đa phương thức (text, hình ảnh, âm thanh, video), tốc độ phản hồi nhanh, tối ưu cho các tác vụ tạo nội dung, tổng hợp thông tin và hội thoại chuyên sâu. Gemini 2.0 Flash nổi bật với khả năng reasoning minh bạch, tích hợp công cụ ngoài và hỗ trợ hơn 100 ngôn ngữ.
    
    \item \textbf{Tavily Web Search:} Công cụ tìm kiếm web thời gian thực dành cho AI, cung cấp API cho phép truy xuất, tổng hợp thông tin mới nhất từ các nguồn uy tín trên internet. Tavily hỗ trợ tích hợp dễ dàng vào workflow của AI agent, đảm bảo dữ liệu luôn được cập nhật, xác thực và phù hợp với nhu cầu của hệ thống.
\end{itemize}

\section{Dataset}

Trong hệ thống Trợ lý Y tế Đa phương thức, bộ dữ liệu VimedAQA được sử dụng làm cơ sở dữ liệu chính cho tác nhân RAG. VimedAQA là bộ dữ liệu y tế tiếng Việt toàn diện, bao gồm nhiều lĩnh vực như bệnh tật, thuốc và các thông tin y tế liên quan khác. Việc sử dụng bộ dữ liệu chuyên biệt và đa dạng này giúp hệ thống có thể truy xuất, hiểu và xử lý các truy vấn y tế bằng tiếng Việt một cách chính xác, đồng thời cung cấp các phản hồi phù hợp với bối cảnh y tế tại Việt Nam.

% Placeholder for image5 - Dataset Overview
\begin{figure}[H]
    \centering
    \fbox{\parbox{0.8\textwidth}{\centering
        \textbf{[CẦN UPLOAD HÌNH ẢNH 5]}\\
        \textit{Tổng quan bộ dữ liệu VimedAQA}\\
        \textit{Mô tả: Biểu đồ thống kê về cấu trúc và phân bổ dữ liệu VimedAQA}
    }}
    \caption{Tổng quan bộ dữ liệu VimedAQA}
    \label{fig:dataset-overview}
\end{figure}

\section{Kết quả \& Đánh giá}

Đề tài không sử dụng các kết quả benchmark để đánh giá mà chỉ tiến hành chạy thực nghiệm mô phỏng trên giao diện web demo:

\begin{itemize}
    \item Giao diện tương tác real-time demo hoạt động ổn định.
    \item Quá trình điều phối, phân tích và sử dụng agent được hiển thị từng bước và cho phép người dùng nhìn thấy kết quả được suy luận từ các nguồn.
\end{itemize}

\subsection{Ưu điểm}

\begin{itemize}
    \item Tính mô-đun cao, dễ mở rộng.
    \item Có khả năng tương tác với dữ liệu multimodal (text + ảnh).
    \item Sử dụng cho domain đặc thù là y tế, có nhiều tiềm năng để phát triển.
    \item Tính thực tế cao với nhu cầu của người dùng, tăng cường hiểu biết và các cách xử lý ban đầu của người dùng với các loại bệnh.
    \item Độ tin cậy cao khi trả lời dựa trên các tài liệu y tế chuyên khoa và các lời tư vấn của bác sĩ trong database hoặc các tổ chức uy tín trên internet.
\end{itemize}

\subsection{Hạn chế}

\begin{itemize}
    \item Dữ liệu y tế cho database còn hạn chế.
    \item Cần thêm một số mô hình AI chuyên biệt để dự đoán các bệnh theo các chuyên ngành đặc biệt.
    \item Cần sử dụng các mô hình LLM có kiến thức về y tế với các agent sử dụng LLM để chẩn đoán bệnh.
    \item Cần sử dụng mô hình embedding tối ưu hơn với nhiều khả năng xử lý được long context tốt hơn, đặc biệt với chuyên ngành y tế có nhiều văn bản dài và phức tạp.
\end{itemize}

Nhìn chung, hướng tiếp cận này phù hợp với thực tiễn phát triển các hệ thống AI y tế hiện đại, tạo nền tảng vững chắc cho việc mở rộng chức năng, nâng cao chất lượng phục vụ và thích ứng linh hoạt với các yêu cầu mới trong lĩnh vực chăm sóc sức khỏe số.

\section{Tiến độ \& Hướng phát triển}

\subsection{Tiến độ hiện tại}

Hoàn thiện các module cốt lõi ở mức cơ bản:

\begin{itemize}
    \item Đã xây dựng và triển khai Decision Agent với chức năng điều phối trung tâm, sử dụng LangGraph để định tuyến yêu cầu đến các agent phù hợp.
    \item Vision Module đã hoàn thành ở mức cơ bản, hỗ trợ nhận diện và phân tích các thông tin hình ảnh, lưu trữ kết quả vào bộ nhớ chia sẻ để duy trì ngữ cảnh.
    \item Text Module đã hoàn thiện:
    \begin{itemize}
        \item Tích hợp RAG Agent để truy xuất thông tin từ cơ sở dữ liệu y tế tiếng Việt (VimedAQA) và Web Search Agent để cập nhật thông tin từ các nguồn uy tín trên internet.
        \item Normal Conversation Agent đã được triển khai, đảm bảo khả năng xử lý các hội thoại thông thường với người dùng.
    \end{itemize}
\end{itemize}

\subsection{Hướng phát triển}

Hướng phát triển trong tương lai của đề tài tập trung vào những điểm sau:

\begin{itemize}
    \item Tăng cường hợp tác và giao tiếp giữa các agent chuyên biệt nhằm nâng cao khả năng xử lý các truy vấn phức tạp, đa chiều và đa nguồn dữ liệu, giúp hệ thống phản ứng linh hoạt và chính xác hơn trong các tình huống đa dạng.

    \item Áp dụng cơ chế học thích nghi (adaptive learning) để hệ thống có thể tự động cải thiện chất lượng phản hồi dựa trên tương tác thực tế với người dùng, từ đó nâng cao độ chính xác và tính phù hợp của các kết quả theo thời gian.
    
    \item Mở rộng khả năng tương tác đa phương thức (multimodal), tích hợp sâu hơn dữ liệu văn bản, hình ảnh, âm thanh và thiết bị y tế đeo tay, giúp hệ thống hỗ trợ chẩn đoán và theo dõi sức khỏe toàn diện hơn.

    \item Phát triển các mô hình AI chuyên biệt cho từng chuyên ngành y tế, như dự đoán bệnh lý chuyên sâu, hỗ trợ chẩn đoán phân biệt, và quản lý bệnh mãn tính cá nhân hóa, nhằm nâng cao giá trị ứng dụng trong thực tiễn.
\end{itemize}

\section{Kết luận}

Dự án đã xây dựng thành công nền tảng hệ thống Trợ lý Y tế Đa phương thức dựa trên kiến trúc multi-agent, đáp ứng tốt các yêu cầu cơ bản về phân tích triệu chứng, truy xuất thông tin y tế và hỗ trợ tư vấn cho người dùng phổ thông tại Việt Nam. 

Hệ thống thể hiện tính mô-đun, khả năng mở rộng và tích hợp đa dạng công nghệ hiện đại, đồng thời đảm bảo độ tin cậy nhờ truy xuất từ các nguồn dữ liệu chuyên ngành và uy tín. Tuy còn một số hạn chế về quy mô dữ liệu và độ sâu chuyên môn của các mô hình AI, kết quả thực nghiệm cho thấy hướng tiếp cận này phù hợp với thực tiễn phát triển các hệ thống AI y tế hiện đại, tạo nền tảng vững chắc cho việc mở rộng chức năng, nâng cao chất lượng phục vụ và thích ứng linh hoạt với các yêu cầu mới trong lĩnh vực chăm sóc sức khỏe số. 

Các hướng phát triển tiếp theo sẽ tập trung vào tối ưu hóa workflow, mở rộng dữ liệu, phát triển các mô hình AI chuyên biệt và tăng cường khả năng tương tác đa phương thức, nhằm nâng cao hiệu quả và giá trị ứng dụng thực tiễn của hệ thống. 